%% kapa: exam=1 round=28 year=2017 period=1 form=A subject=economics subject_ko=경제학원론 seq=13 number=53
%% source: Q-Net 감정평가사 기출문제 (cst00309) artlSeq=5199603 file=extracted/1008/5199603/2017년 제28회 감정평가사 1차시험 문제지/1교시 A형.pdf
%% conversion: kordoc 4.24.51 -> tools/gen1.py (auto); review=proofread-sonnet; flags=none
\begin{question}{53}
복점(duopoly)시장에서 기업 A와 B는 각각 1, 2, 3의 생산량 결정 전략을 갖고 있다. 성과보수행렬(payoff matrix)이 다음과 같을 때 내쉬균형은? (단, 게임은 일회성이며, 보수행렬 내 괄호 안 왼쪽은 A, 오른쪽은 B의 보수이다.)
\begin{center}
\renewcommand{\arraystretch}{1.15}
\begin{tabular}{|c|c|c|c|c|}\hline
\multicolumn{2}{|c|}{} & \multicolumn{3}{c|}{B} \\ \cline{3-5}
\multicolumn{2}{|c|}{} & 전략 1 & 전략 2 & 전략 3 \\ \hline
\multirow{3}{*}{A} & 전략 1 & (7, 7) & (5, 8) & (4, 9) \\ \cline{2-5}
 & 전략 2 & (8, 5) & (6, 6) & (3, 4) \\ \cline{2-5}
 & 전략 3 & (9, 4) & (4, 3) & (0, 0) \\ \hline
\end{tabular}
\end{center}
\begin{choices}
  \item (7, 7), (6, 6), (0, 0)
  \item (7, 7), (5, 8), (9, 4)
  \item (8, 5), (6, 6), (3, 4)
  \item (9, 4), (5, 8), (0, 0)
  \item (9, 4), (6, 6), (4, 9)
\end{choices}
\end{question}
\answer{5}
